\chapter*{МЭРГЭЖЛИЙН НЭР ТОМЬЁОНЫ ТАЙЛБАР}
\addcontentsline{toc}{chapter}{Мэргэжлийн нэр томьёоны тайлбар}

\begin{table}[H]
  \centering
  \caption{Мэргэжлийн нэр томьёо}
  \label{tab:glossary}
  \begin{tabularx}{\textwidth}{>{\raggedright\arraybackslash}p{4.5cm} X}
    \toprule
    Нэр томьёо & Товч тайлбар \\
    \midrule
    SISi системийн ухаалаг агент (SISi Intelligent Agent) & SISI порталд суурилсан хиймэл оюуны туслах систем. Хэрэглэгчийн байгалийн хэлний хүсэлтийг бодит системийн үйлдэл болгон хувиргах зориулалттай. \\
    \addlinespace
    Model Context Protocol (MCP) & Хиймэл оюуны загварт гадаад хэрэгсэл, өгөгдлийн эх үүсвэрийг стандартчилсан аргаар холбох нээлттэй протокол. tool, resource, prompt гэсэн үндсэн нэгжүүд ашиглана. \\
    \addlinespace
    AI Agent (Хиймэл оюуны агент) & Хэлний загвар, орчны өгөгдөл, хэрэгслүүд, үйлдлийн урсгалыг нэгтгэсэн ухаалаг систем. Зөвхөн хариу үүсгэх бус, бодит үйлдэл гүйцэтгэх чадвартай. \\
    \addlinespace
    Tool-calling & Хиймэл оюун бодит системтэй харилцахдаа хэрэгслүүдийг дуудаж, өгөгдөл авах, үйлдэл гүйцэтгэх механизм. \\
    \addlinespace
    MCP Server & Model Context Protocol-ийн стандартад нийцсэн сервер. Гаднын системүүдийг (SISI портал, календарь, баримт бичиг гэх мэт) хиймэл оюунтай холбох давхарга. \\
    \addlinespace
    Session Bootstrap & Хэрэглэгчийг дахин нэвтрүүлэхгүйгээр одоогийн access token, cookie ашиглан сесс үүсгэх механизм. \\
    \addlinespace
    Access Token & Хэрэглэгчийн системд нэвтэрснийг баталгаажуулж, API хүсэлтийг зөвшөөрөх нууцлалын түлхүүр. \\
    \addlinespace
    Read-only Tool & Өгөгдөл унших, харах зориулалттай MCP хэрэгсэл. Mutation үйлдэл хийх боломжгүй. \\
    \addlinespace
    Mutation & Өгөгдөл бичих, өөрчлөх, устгах зэрэг системд өөрчлөлт оруулах үйлдэл. Баталгаажуулалт шаарддаг. \\
    \addlinespace
    Chrome Extension & Вэб хөтөчид нэмэлт функц нэмэх програм. Энэхүү төслөөр порталын сессийг авах, хэрэглэгчийн интерфейс харуулахад ашиглана. \\
    \addlinespace
    Tool Orchestration & Олон хэрэгслийг дарааллан эсвэл параллелаар дуудаж, үр дүнг нэгтгэх үйл явц. \\
    \addlinespace
    Domain & Системийн тодорлхой чиглэл (хуваарь, дүн, санхүү гэх мэт). \\
    \addlinespace
    Prompt & Хиймэл оюуны загварт өгөх заавар, асуулт, эсвэл контекст. \\
    \addlinespace
    Orchestration Layer & AI сервер дээр хэрэгслүүдийг дуудах, хариулт үүсгэх үйл явцыг удирдах модуль. \\
    \addlinespace
    Audit Trail & Mutation үйлдлүүдийн түүх, лог бичлэг. Аюулгүй байдал, хяналтын зориулалттай. \\
    \addlinespace
    Confirmation Flow & Эрсдэлтэй үйлдэл хийхээс өмнө хэрэглэгчээс дахин баталгаа авах механизм. \\
    \addlinespace
    Workflow & Олон алхамт үйлдлүүдийн дараалал. Жишээ нь: Excel файлаас дүн оруулах. \\
    \addlinespace
    API (Application Programming Interface) & Системүүд хоорондоо өгөгдөл солилцох протокол. \\
    \addlinespace
    Elysia & TypeScript хэл дээр бичигдсэн хөнгөн жинтэй вэб фреймворк. \\
    \addlinespace
    React & Хэрэглэгчийн интерфейс бүтээхэд ашиглах frontend сан. \\
    \addlinespace
    Bun & JavaScript/TypeScript код ажиллуулах шинэ runtime. Хурдан, хөнгөн. \\
    \addlinespace
    Turborepo & Monorepo удирдах хэрэгсэл. Олон апп, package-ийг нэг кодын сангаас зохицуулна. \\
    \addlinespace
    Drizzle ORM & TypeScript-д зориулсан өгөгдлийн сантай харилцах сан. \\
    \bottomrule
  \end{tabularx}
\end{table}

\vspace{1cm}
\noindent\textbf{Тэмдэглэл:} Энэхүү тайландаа анх удаа гарч буй мэргэжлийн нэр томьёог товч тайлбартай хамт оруулсан. Давтагдах үед хэвээрээ англи нэр томьёог ашиглана.
